\chapter{関連する研究と事例}\label{ch:related}

本報告書で扱った現象は、AMBAC という名前が付く前から、生物学・スポーツ科学・宇宙工学・ロボット工学のさまざまな分野で研究されてきた。
本章では、本報告書の結果と対応づけながら、代表的な研究と事例を紹介する。

\section{落下する猫と宇宙飛行士}

\paragraph{落下する猫}
背中を下にして落とされた猫が、空中で体をひねって足から着地する現象は、角運動量がゼロのまま向きを変える典型例である。
Kane と Scher は 1969 年に、猫の体を 2 つの円柱がつながったものとみなして、この現象を力学的に説明した \cite{R7}。
Montgomery はこれをゲージ理論（接続と曲率）の言葉で整理した \cite{R4}。本報告書の第~\ref{ch:theory}章、第~\ref{ch:planar}章は、その平面版にあたる。

\paragraph{宇宙飛行士の姿勢変更}
Kane と Scher は 1970 年に、宇宙飛行士が手足の動きで体の向きを変える方法を解析した \cite{R8}。
近年では、宇宙飛行士の大西卓哉さんが国際宇宙ステーションで脚を使って姿勢を制御し、AMBAC の名を挙げたことが話題になった \cite{R43,R44}。
第~\ref{ch:legs}章の「宇宙では脚が最も効く重い手足」という結果は、この実例と整合する。

\section{スポーツ：体操・飛び込みのひねり}

体操や飛び込みの選手は、空中で腕などを非対称に動かすことで、宙返りにひねりを加える。
Yeadon は、宙返りの最中に腕・胸・腰を矢状面に対して非対称に動かすと、体の軸が傾き、それがひねりを生むことを示した \cite{R46}。
宙返りがない場合は、1 回のひねりに何回も手足を往復させる必要があるとされる。
これは第~\ref{ch:3d}章の 3 次元の AMBAC（円錐運動による姿勢変更、回転の非可換性）が、実際の人の運動で使われている例である。

\section{動物とロボットの尻尾}

Libby らは、トカゲが跳躍するときに尻尾を振って体の角運動量を尻尾に移し、空中で体の前後の傾き（ピッチ）を安定させていることを示し、
尻尾を持つ小型ロボットで同じ制御を再現した \cite{R45}。
尻尾は、本報告書の言葉では「重い手足を 1 本だけ持つ AMBAC」であり、空中の短い時間で可逆な首振り（第~\ref{ch:planar}章）を使って姿勢を立て直している。
第~\ref{ch:legs}章で、脚（重い手足）を重くするほど AMBAC が強くなったことと対応する。

\section{宇宙ロボット}

宇宙ステーションや人工衛星に取り付けたロボットアームを動かすと、反作用で母機の姿勢が乱れる。
この問題は 1980〜90 年代から研究され、アームの動きと母機の姿勢の関係を表す一般化ヤコビ行列 \cite{R14}、
母機の姿勢を乱さないアームの動かし方（反動零空間制御）\cite{R15}、
アームのループ運動で母機の姿勢を変える経路計画 \cite{R12}、制御則の性質の解析 \cite{R13} などが提案されている。
本報告書の AMBAC の基本式 $\dot\theta=-\Gamma\cdot\dot q$ は、一般化ヤコビ行列の姿勢成分にあたる。

\section{ヒューマノイドのバランス}

二足歩行ロボットのバランス制御では、ZMP の考え方 \cite{R28} が広く使われてきた。
人の立位姿勢の研究では、外乱に対して足首を使う「足首の戦略」と、上体を使う「腰の戦略」が知られている \cite{R31}。
第~\ref{ch:legs}章、第~\ref{ch:steps}章で、二本脚が荷物を運んだり段差を上ったりするのに上体の前後傾が不可欠だったことは、腰の戦略にあたる。

アクチュエータについては、代表的なヒューマノイドである Boston Dynamics の Atlas が、2024 年に油圧式から電動式に切り替わった \cite{R47}。
人間程度の大きさでは電動でも十分な性能が得られる、という第~\ref{ch:conclusion}章の議論と整合する。

\section{学習による動作の獲得}

モーションキャプチャーのデータを手本に、物理シミュレーション上で強化学習により動作を獲得する手法 \cite{R24}、
周期的な軌道生成器を学習で調整する手法 \cite{R21}、シミュレーションのパラメータを乱数で変えて実機への移行を容易にする手法 \cite{R23} が提案されている。
第~\ref{ch:ai}章の結果（構造を持たない方策の失敗、軌道生成器型の成功、ドメインランダム化による頑健さ）は、これらの研究の考え方を AMBAC という具体例で確かめたものと言える。

\section{本報告書の位置づけ}

\begin{longtable}{p{0.24\linewidth}p{0.32\linewidth}p{0.36\linewidth}}
\toprule
分野 & 代表的な研究・事例 & 本報告書との対応 \\ \midrule \endhead
生物・幾何学 & 落下する猫 \cite{R7,R4} & 幾何学的位相（第~\ref{ch:theory}, \ref{ch:planar}章） \\
宇宙飛行士 & 手足による姿勢変更 \cite{R8,R43,R44} & 宇宙での脚の価値（第~\ref{ch:legs}章） \\
スポーツ科学 & 宙返りのひねり \cite{R46} & 3 次元の非可換性（第~\ref{ch:3d}章） \\
動物・ロボット & 尻尾による姿勢制御 \cite{R45} & 重い手足 1 本の首振り（第~\ref{ch:planar}, \ref{ch:legs}章） \\
宇宙ロボット & 一般化ヤコビ行列、反動零空間 \cite{R14,R15} & AMBAC の基本式（第~\ref{ch:theory}, \ref{ch:rw}章） \\
ヒューマノイド & ZMP、腰の戦略 \cite{R28,R31} & 転倒・段差（第~\ref{ch:tank}, \ref{ch:legs}, \ref{ch:steps}章） \\
学習 & 模倣学習、PMTG、ドメインランダム化 \cite{R24,R21,R23} & 学習の成否と頑健さ（第~\ref{ch:ai}章） \\
\bottomrule
\end{longtable}

本報告書の新しさは、これらの別々の分野の現象を、フィクションの AMBAC を軸に 1 つの枠組み（接続・曲率・角運動量の容量・接触の限界・相似則）でつなぎ、
形状の優位が用途と大きさで入れ替わることを定量的に示した点にある。
